\documentclass[12pt]{article}
\usepackage{amsmath, amssymb, amsthm, mathtools, hyperref}


\newcommand{\A}{\mathcal{A}}
\newcommand{\B}{\mathcal{B}}
\newcommand{\N}{\mathbb{N}}
\newcommand{\AUT}{\mathrm{AUT}}

\newtheorem{definition}{Definition}
\newtheorem{question}{Question}
\newtheorem*{theorem*}{Theorem}

\begin{document}
\title{Submission D solution to Problem 1}
\date{}
\maketitle

We consider infinite countable first-order structures, their automorphism groups and the descriptive aspects of their representations. A representation $\A$ of such a structure is an instance of its isomorphism type in which the universe or domain of the structure is the set of natural numbers $\N$ and the constants, functions and relations of the structure are elements of $\N$, multi-variable functions from $\N$ to $\N$ and multi-place relations on $\N$, respectively.   Similarly, an automorphism of $\A$ or an isomorphism from $\A$ to $\B$ is a bijective function from $\N$ to $\N$ that preserves the evaluation of all constants, functions and relations.  We let $\AUT(\A)$ denote the automorphism group of $\A$.

\begin{definition}\label{1.1}
  \begin{itemize}
  \item     $\A$ is \emph{computably $\AUT$-countable} if every automorphism of $\A$ is computable relative to $\A$.  
  \item  $\A$ is \emph{computably $\AUT$-countable on a cone} if there is a $C\subset\N$ such that whenever $\B$ is isomorphic to $\A$, every automorphism of $\B$ is computable relative to $\B\oplus C$.
  \end{itemize}
\end{definition}


\begin{definition}
  A function $f:\A\to\A$ is \emph{$\Sigma^{in}_1$-definable} in $\A$ relative to parameters
  $\bar{a}=(a_1,\dots,a_n)$ from $\A$ if there is a family of existential formulas $(\phi_i)_{i\in\N}$ such that for all $x$ and $y$ in $\A$, $f(x)=y$ if and only if there is an $i\in\N$ such that $\A\models\phi_i(\bar{a},x,y)$.
\end{definition}

\begin{question}\label{sec:question}\label{1.2}
  Is there a countable computably-represented structure $\A$ such that $\A$ is computably $\AUT$-countable on a cone and such that for any finite  set of parameters $\bar{a}$ in $\A$ there is an automorphism $\pi$ of $\A$ such that $\pi$ is not $\Sigma^{in}_1$-definable in $\A$ relative to the parameters $\bar{a}\cup \pi(\bar{a}),$ where $\pi(\bar{a})$ is the image of $\bar{a}$ under $\pi$.
\end{question}

\section*{Solution}

\begin{proof}
We will prove that the answer to the question is negative: there is no such structure. We proceed by contradiction. Assume that there exists a countable computably-represented structure $\mathcal{A}$ satisfying the two conditions in the problem.

\textbf{Step 1: The automorphism group of $\mathcal{A}$ is countable.}

By the first condition, $\mathcal{A}$ is computably $\AUT$-countable on a cone. Thus, there exists an oracle $C \subset \mathbb{N}$ such that for every structure $\mathcal{B}$ isomorphic to $\mathcal{A}$, every automorphism of $\mathcal{B}$ is computable relative to $\mathcal{B} \oplus C$. 
Applying this to $\mathcal{B} = \mathcal{A}$, we have that every automorphism of $\mathcal{A}$ is computable relative to the oracle $\mathcal{A} \oplus C$. Because there are only countably many Turing machines, there are only countably many functions that are computable relative to any fixed oracle. Therefore, the set of all automorphisms of $\mathcal{A}$ is countable, which means the group $\AUT(\mathcal{A})$ is countable.

\textbf{Step 2: The group $\AUT(\mathcal{A})$ contains an isolated point.}

The group $\AUT(\mathcal{A})$ is a closed subgroup of the symmetric group $S(\mathbb{N})$ equipped with the topology of pointwise convergence. Therefore, $\AUT(\mathcal{A})$ is a Polish space (a separable and completely metrizable topological space). 
By the Baire Category Theorem, any non-empty countable Polish space must contain at least one isolated point. Because $\AUT(\mathcal{A})$ is a topological group, its topology is homogeneous, meaning that if one point is isolated, then all points are isolated. In particular, the identity element $\mathrm{id} \in \AUT(\mathcal{A})$ is an isolated point.

\textbf{Step 3: There exists a finite tuple with a trivial pointwise stabilizer.}

In the topology of pointwise convergence on $\AUT(\mathcal{A})$, the basic open neighborhoods of the identity $\mathrm{id}$ are the pointwise stabilizers of finite tuples of elements from $\mathbb{N}$. 
Because $\mathrm{id}$ is an isolated point, the singleton set $\{\mathrm{id}\}$ is open. Thus, there exists a basic open neighborhood of $\mathrm{id}$ that is exactly $\{\mathrm{id}\}$. This implies there exists a finite tuple $\bar{c}$ in $\mathcal{A}$ such that if $\sigma \in \AUT(\mathcal{A})$ and $\sigma(c_i) = c_i$ for all elements $c_i$ in $\bar{c}$, then $\sigma = \mathrm{id}$.

\textbf{Step 4: Fixing the parameters and the automorphism $\pi$.}

We choose our finite set of parameters to be the tuple $\bar{a} = \bar{c}$. By the second condition of the problem, there exists an automorphism $\pi$ of $\mathcal{A}$ such that $\pi$ is not $\Sigma^{in}_1$-definable in $\mathcal{A}$ relative to the parameters $\bar{c} \cup \pi(\bar{c})$. 
Let $\bar{d} = \pi(\bar{c})$. 
We define the expanded structure $\mathcal{A}^* = (\mathcal{A}, \bar{c}, \bar{d})$, which is the structure $\mathcal{A}$ together with new constant symbols evaluating to the elements of $\bar{c}$ and $\bar{d}$. 

\textbf{Step 5: Rigidity of the expanded structure $\mathcal{A}^*$.}

Let $\sigma$ be an automorphism of the expanded structure $\mathcal{A}^*$. By definition, $\sigma$ is an automorphism of $\mathcal{A}$ that fixes the interpretations of the new constants, meaning $\sigma(\bar{c}) = \bar{c}$ and $\sigma(\bar{d}) = \bar{d}$. 
Because $\sigma(\bar{c}) = \bar{c}$, Step 3 implies that $\sigma = \mathrm{id}$. 
Therefore, the only automorphism of $\mathcal{A}^*$ is the identity. Consequently, every relation on the domain of $\mathcal{A}^*$ is invariant under all automorphisms of $\mathcal{A}^*$. 
In particular, the binary relation $R = \{(x, y) \in \mathbb{N}^2 : \pi(x) = y\}$ (the graph of $\pi$) is invariant under all automorphisms of $\mathcal{A}^*$.

\textbf{Step 6: Application of the Ash--Knight--Manasse--Slaman and Chisholm Theorem.}

We rely on the following theorem from computable structure theory:
\begin{theorem*}[Ash--Knight--Manasse--Slaman and Chisholm]
Let $M$ be a countable first-order structure, and let $R$ be a relation on the domain of $M$. Suppose that $R$ is invariant under all automorphisms of $M$. Suppose there exists an oracle $C \subset \mathbb{N}$ such that for every structure $N$ with domain $\mathbb{N}$ and every isomorphism $f: M \to N$, the relation $f(R)$ is computable relative to $N \oplus C$. Then $R$ is definable in $M$ by an $L_{\omega_1\omega}$ $\Sigma_1$ formula without parameters (meaning there is a countable family of parameter-free existential formulas in the language of $M$ whose disjunction defines $R$).
\end{theorem*}

We verify the hypotheses of this theorem for $M = \mathcal{A}^*$ and $R = \operatorname{graph}(\pi)$:
\begin{enumerate}
\item $M = \mathcal{A}^*$ is a countable first-order structure, as $\mathcal{A}$ has domain $\mathbb{N}$ and we added finitely many constants.
\item $R$ is a relation on the domain of $M$.
\item $R$ is invariant under all automorphisms of $M$, which was established in Step 5.
\item Let $C$ be the fixed oracle from Step 1. Let $N$ be any structure with domain $\mathbb{N}$ and let $f: \mathcal{A}^* \to N$ be an isomorphism. The structure $N$ is an expansion of a structure $\mathcal{B}$ by constants $f(\bar{c})$ and $f(\bar{d})$, where $\mathcal{B}$ is a structure in the language of $\mathcal{A}$. The map $f$ is an isomorphism from $\mathcal{A}$ to $\mathcal{B}$. 

The relation $f(R)$ is $\{ (f(x), f(y)) : \pi(x) = y \}$, which is exactly the graph of the function $f \circ \pi \circ f^{-1}$. Because $\pi$ is an automorphism of $\mathcal{A}$ and $f: \mathcal{A} \to \mathcal{B}$ is an isomorphism, $f \circ \pi \circ f^{-1}$ is an automorphism of $\mathcal{B}$. 

By the definition of $C$, every automorphism of $\mathcal{B}$ is computable relative to $\mathcal{B} \oplus C$. Therefore, $f \circ \pi \circ f^{-1}$ is computable relative to $\mathcal{B} \oplus C$. Because the atomic diagram of $N$ (viewed as a real in $2^\omega$) consists of the atomic diagram of $\mathcal{B}$ and the evaluations of the finitely many constants, the oracle $\mathcal{B} \oplus C$ is Turing equivalent to $N \oplus C$. Thus, $f(R)$ is computable relative to $N \oplus C$. This verifies the fourth hypothesis.
\end{enumerate}

\textbf{Step 7: Reaching a contradiction.}

By the theorem, $R$ is definable in $\mathcal{A}^*$ by a countable disjunction of parameter-free existential formulas. This means there is a family of existential formulas $(\phi_i(x,y))_{i \in \mathbb{N}}$ in the language of $\mathcal{A}^*$ such that for all $x$ and $y$ in $\mathbb{N}$, $R(x, y)$ holds if and only if there is an $i \in \mathbb{N}$ such that $\mathcal{A}^* \models \phi_i(x, y)$.

The language of $\mathcal{A}^*$ is the language of $\mathcal{A}$ expanded by constant symbols for $\bar{c}$ and $\bar{d}$. By replacing these constant symbols with variables, we obtain a family of existential formulas $(\psi_i(\bar{u}, \bar{v}, x, y))_{i \in \mathbb{N}}$ in the language of $\mathcal{A}$ such that for all $x$ and $y$ in $\mathbb{N}$, $\pi(x) = y$ if and only if there is an $i \in \mathbb{N}$ such that $\mathcal{A} \models \psi_i(\bar{c}, \bar{d}, x, y)$.

This matches the exact definition given in the problem statement for the function $\pi$ being $\Sigma^{in}_1$-definable in $\mathcal{A}$ relative to the parameters $\bar{c} \cup \bar{d}$. 
Because $\bar{d} = \pi(\bar{c})$, we conclude that $\pi$ is $\Sigma^{in}_1$-definable in $\mathcal{A}$ relative to the parameters $\bar{c} \cup \pi(\bar{c})$. 
This directly contradicts the initial choice of $\pi$ from Step 4. 

\textbf{Conclusion:}

Because the assumption that such a structure $\mathcal{A}$ exists leads to a logical contradiction, we conclude that no such countable computably-represented structure exists. The answer to Question \ref{1.2} is no.
\end{proof}

\end{document}